% !TeX root = ../动手学强化学习笔记.tex
% 由主文档通过 \input 载入；请编译主文档。
\setcounter{chapter}{11} % 保留原书章号。
\chapter{价值函数近似}
\label{ch:rl-approximation}

表格方法在每个状态保存独立信息，连续或高维状态下很难充分访问。函数近似（function approximation）用共享参数把已见经验推广到未见状态；与此同时，一次更新会影响多个预测，表格算法的稳定性论证因此不再直接适用。

\section{特征、表示与误差目标}

令特征向量 $\phi(s)\in\mathbb R^p$、参数 $w\in\mathbb R^p$，线性价值为 $V_w(s)=\phi(s)^\top w$。把全部状态特征排列为矩阵 $\Phi\in\mathbb R^{n\times p}$，则预测向量为 $\Phi w$。状态聚合使用组指示特征，神经网络则同时学习非线性表示与输出映射。

给定状态权重分布 $d$，定义 $D=\operatorname{diag}(d)$ 和平方价值误差
\[
L_{\rm VE}(w)=\frac12\sum_s d(s)(V_w(s)-v_\pi(s))^2.
\]
不同 $d$ 会产生不同最佳近似；数据很少的状态可能几乎不影响损失。真实价值未知时不能直接计算该损失，需要回报样本或自举目标。表示能力、样本误差与优化误差应分开讨论。

\section{蒙特卡洛梯度与半梯度 TD}

用固定策略回报 $G$ 作为监督目标，单样本梯度更新为 $w\leftarrow w+\alpha(G-V_w(S))\nabla_wV_w(S)$。在条件无偏和可交换微分期望等条件下，它对应平方价值误差的随机梯度。

半梯度（semi-gradient）TD 则把 $Y=R+\gamma V_w(S')$ 在本次更新中当作常数：
\begin{equation}
w\leftarrow w+\alpha\bigl[R+\gamma V_w(S')-V_w(S)\bigr]\nabla_wV_w(S).
\label{eq:rl-semi-gradient}
\end{equation}
它没有对目标中的 $V_w(S')$ 求导，因此一般不是某个简单平方 Bellman 残差的完整梯度。对平方期望贝尔曼残差求梯度还涉及条件期望乘积，单个后继样本一般不能直接给出无偏乘积估计，这称为双采样困难。

\begin{example}[共享参数的影响]
取 $\phi(s_1)=1,\phi(s_2)=2,w=1$，故两个估计为 $1,2$。观察 $s_1\to s_2$、奖励零，$\gamma=0.9,\alpha=0.1$，半梯度更新得到 $w=1.08$，两个估计同时变为 $1.08,2.16$。若固定数据只反复提供这一转移，参数递推为 $w_{k+1}=1.08w_k$，将发散；这个缺乏后继覆盖的例子说明局部自举误差可能通过共享参数放大。
\end{example}

\section{投影方程与稳定性}

在有限状态、固定策略下，假设 $d(s)>0$ 是 $P_\pi$ 的平稳分布，$\Phi$ 满列秩。以 $\|x\|_D^2=x^\top Dx$ 定义加权范数，令 $\Pi$ 为到特征张成空间的正交投影。平稳性与 Jensen 不等式给出 $\|P_\pi x\|_D\leq\|x\|_D$，而正交投影不扩大距离，因此 $\Pi T_\pi$ 是压缩映射。线性 TD 的平均方程为
\[
\Phi w=\Pi(r_\pi+\gamma P_\pi\Phi w),\qquad
Aw=b,
\]
其中 $A=\Phi^\top D(I-\gamma P_\pi)\Phi$、$b=\Phi^\top Dr_\pi$。对非零 $w$，$w^\top Aw\geq(1-\gamma)\|\Phi w\|_D^2>0$，所以 $A$ 可逆。这里的固定点一般不同于最佳价值投影 $\Pi v_\pi$。

这解释了同策略线性近似的稳定结构，但随机迭代收敛仍需遍历采样、步长与噪声等条件。离策略采样的权重分布未必是 $P_\pi$ 的平稳分布，上述范数论证可能失效；神经网络的表示空间又非固定线性子空间。函数近似、自举与离策略同时出现时需特别检验稳定性，而不是认为组合必然成功或必然发散。

\section{神经网络近似的训练}

网络将状态映射为价值或一组动作价值，优化器只根据输入批次和目标调整参数。输入尺度、奖励尺度、初始化与学习率会改变梯度大小；目标网络、梯度裁剪等措施可以影响数值稳定性，但不能替代正确的目标与数据覆盖。训练损失下降可能只是网络更好地拟合当前错误目标，策略回报必须独立评价。

实践中应先验证小表格问题，再换成线性特征，最后使用网络。每次只增加一种复杂性，才能辨认错误来自环境、回报、共享表示还是优化过程。

\section{习题与实验}
\begin{enumerate}
\item 两状态真实价值为 $1,3$，近似被限制为同一常数。分别在均匀权重和权重 $(0.9,0.1)$ 下求最佳近似。
\item 推导线性 MC 与线性半梯度 TD 更新，指出完整样本平方 TD 误差梯度中被半梯度忽略的项。
\item 在小 MRP 上分别计算最佳投影和 TD 固定点；改变特征与采样分布，比较解析解和迭代结果，并检查矩阵是否可逆。
\end{enumerate}





